\subsection{章末確認問題}
\begin{itemize}
\item ソフトウェア保守は、なぜ単なるバグ修正ではないと言えるか。
\item 是正保守、予防保守、適応保守、追加保守、完全化保守の違いを説明せよ。
\item 緊急保守は、是正保守とどのように違うか。
\item 保守担当者の「限られた理解」は、なぜ保守費用を増やすのか。
\item 回帰テストが保守で重要になる理由を説明せよ。
\item 影響分析では、どのような情報を確認すべきか。
\item 保守性が低いソフトウェアは、どのように技術的負債を生むか。
\item 開発チームと保守チームを同一にする場合と分離する場合の利点と注意点を説明せよ。
\item 継続的統合、継続的デリバリ、継続的テスト、継続的デプロイは、保守にどう役立つか。
\item 保守ツールの中から三つを選び、何を支援するか説明せよ。
\end{itemize}
\begin{pointbox}{解答の方向性}保守は、運用中のソフトウェアを支援し、変更し、価値を保つ活動である。分類問題では、欠陥修正か、潜在欠陥の予防か、環境適応か、機能追加か、品質改善かを見分ける。影響分析では、変更対象、依存関係、テスト、性能、安全性、セキュリティ、費用、リスクを確認する。\end{pointbox}
